\section{Conclusiones}

El análisis estadístico e inferencial realizado sobre el dataset \textbf{Diabetic
Retinopathy Debrecen} permitió caracterizar la naturaleza distribucional de los
biomarcadores de microaneurismas, validar la existencia de diferencias sistemáticas
entre grupos y evaluar la capacidad predictiva de un modelo lineal de referencia.
Los resultados obtenidos responden de manera coherente a la pregunta clínica central
del estudio y establecen una base metodológica sólida para el desarrollo de
clasificadores de mayor complejidad.

\subsection{Síntesis de Hallazgos Principales}

Se confirmó la \textbf{no normalidad} de todas las variables de microaneurismas mediante
la prueba de Shapiro-Wilk, con p-valores del orden de $10^{-20}$ que justifican
inequívocamente el uso de estadística no paramétrica en el análisis inferencial. Esta
característica distribucional, junto con la asimetría positiva documentada, es
clínicamente coherente: en una población que incluye tanto pacientes sanos como
afectados, la mayoría de los individuos presentan recuentos bajos de lesiones, mientras
que una minoría con enfermedad avanzada genera la cola derecha de la distribución.

La validación de valores atípicos mediante el método de Tukey reforzó esta interpretación:
la práctica totalidad de los outliers identificados corresponde a la Clase 1, confirmando
que los recuentos extremos de microaneurismas son manifestaciones clínicas genuinas de
estadios severos de retinopatía diabética y no errores de medición. La única excepción
(índice 14, Clase 0, \texttt{ma6}) es explicable por la naturaleza probabilística del
sistema de detección al operar con umbral $\alpha = 1{,}0$.

\subsection{Discusión de Resultados y Desempeño del Modelo}

La prueba de \textbf{Mann-Whitney U} demostró diferencias estadísticamente significativas
entre grupos en todos los niveles de confianza (p $<$ 0,05 en todos los casos), siendo
\texttt{ma1} la variable con mayor potencia discriminatoria ($p = 1{,}25 \times 10^{-21}$).
El gradiente de significancia observado desde \texttt{ma1} hasta \texttt{ma6} evidencia
que los detectores de menor umbral capturan la señal clínica más nítidamente, respaldando
su uso prioritario en sistemas de cribado temprano. El \textbf{Test de Levene} añadió
una dimensión clínica relevante a este análisis: los grupos difieren no solo en tendencia
central, sino también en dispersión, lo que refleja la heterogeneidad inherente al
espectro de severidad de la retinopatía diabética.

El modelo de \textbf{Regresión Logística}, concebido explícitamente como clasificador
de referencia (\textit{baseline}), confirmó la significancia individual de \texttt{ma1}
y \texttt{exudate1} como predictores primarios (p = 0,000 en ambos casos), mientras que
los rasgos anatómicos estructurales no mostraron aporte predictivo lineal. No obstante,
el desempeño global del modelo evidenció limitaciones críticas: una sensibilidad del
61,7\%, una especificidad del 56,3\% y un AUC de 0,65. Estos valores confirman que,
si bien los biomarcadores seleccionados contienen información diagnóstica real, la
relación entre predictores y diagnóstico presenta una complejidad no lineal que no puede
ser adecuadamente capturada mediante una función de decisión lineal. El pseudo-$R^{2}$
de McFadden de 0,087 cuantifica esta limitación: el modelo explica menos del 9\% de la
varianza en el diagnóstico, subrayando que el problema de clasificación involucra
interacciones entre variables y patrones de alta dimensionalidad que escapan al alcance
de la regresión logística estándar.

\subsection{Limitaciones y Proyecciones Futuras}

Las principales limitaciones de este análisis son de dos órdenes. En primer lugar, de
naturaleza metodológica: la ausencia de medidas de tamaño de efecto en la comparación
de grupos impide cuantificar la magnitud práctica de las diferencias detectadas más
allá de su significancia estadística; su incorporación mediante métricas como la
correlación de rango biserial o el estadístico $\hat{A}_{12}$ fortalecería la
interpretabilidad clínica de los resultados. En segundo lugar, de naturaleza estructural:
la alta colinealidad entre las variables MA ($r \approx 1{,}00$), aunque gestionada
mediante la selección de \texttt{ma1} como representante del grupo, sugiere que el
espacio de características podría beneficiarse de una reducción de dimensionalidad
formal, como el Análisis de Componentes Principales (PCA), preservando la varianza
diagnóstica con menor redundancia.

Como proyección futura, los resultados de este estudio señalan de forma inequívoca la
necesidad de algoritmos de inteligencia computacional no lineales. Métodos como
\textit{Random Forests}, \textit{Gradient Boosting} o redes neuronales convolucionales
aplicadas directamente sobre imágenes de fondo de ojo están en capacidad de explotar la
alta colinealidad detectada como información complementaria, modelar interacciones no
lineales entre lesiones y rasgos anatómicos, y alcanzar los niveles de sensibilidad
($>$90\%) y AUC ($>$0,95) exigidos para sistemas de cribado automatizado con
aplicabilidad clínica real, tal como documentan \textcite{casanova2014application}. El
modelo logístico desarrollado en este trabajo proporciona el umbral de comparación
cuantitativo indispensable para evaluar el beneficio marginal de esos enfoques más
sofisticados.
